\justifying
\NSFCBodyText

\subsubsection{研究背景与意义}

新疆口岸连接我国与中亚多条跨境运输通道，货物经历长距离运输、频繁装卸和边境仓储等连续环节。运输途中持续振动、异常倾斜、非法开启及包装破损会改变货物完整性；进入仓储后，温湿度偏移、烟雾、异常门禁和视频事件又可能形成新的环境与安防风险。物流物联网研究已经表明，多源感知、边缘计算和混合网络有助于提高运输与仓储过程的可见性，但既有成果多集中于系统架构、业务效率或单场景原型，复杂网络与资源约束下的可靠异常判断仍缺少统一的可重复验证\cite{internetofthings2019,iotbasedsupply2023,applicationsofthe2022,maritimeterminalscargo2023,hybridiotnetwork2025}。

这一场景同时受到三类约束。其一，振动、倾斜、电子封志、光感和包装图像具有不同采样频率、噪声机制和故障模式，网络中断还会造成异步、缺失和延迟，使真实货物异常与观测质量下降相互混淆。其二，2G/3G 等弱网中的带宽、送达时延、端侧算力和电池能量共同受限，全量采集、完整推理和持续上传不可兼得，离线缓存、断点续传、压缩、边缘计算与低功耗管理必须围绕风险任务协同。其三，温湿度缓变与烟雾、门禁、视频突发具有不同时标，独立阈值难以识别跨源共现和风险累积，也难以在控制误报的同时获得稳定提前量。

因此，本项目拟研究“弱网环境下跨境物流多模态智能感知与边缘协同方法”，不是简单增加传感器或搭建综合物流平台，而是回答三个相互衔接的认知问题：不完备观测下多模态互补信息如何形成可信的货物状态表征；通信、计算和能量如何依据风险信息价值进行分配；不同时标的仓储多源证据如何形成可校准的复合风险预警。三条主线分别对应运输感知、弱网边缘协同和仓储预警，并共享事件风险、信息可靠性与资源约束三个核心变量，可为边远地区移动感知、冷链和跨境物流等场景提供可迁移的方法依据。

\subsubsection{国内外研究现状及发展动态分析}

\textbf{（1）货物多模态感知正在由完整输入下的融合转向不完备观测下的可靠表征。}多模态学习已形成表示、对齐、融合和协同学习的基本框架\cite{multimodalmachinelearning2018,multimodallearningwith2023}。针对测试阶段模态缺失，EmbraceNet 等方法通过训练期随机丢弃增强鲁棒性，多模态 Transformer 的系统评价进一步揭示了缺失输入造成的显著退化\cite{embracenetarobust2019,aremultimodaltransformers2022}。更值得注意的是，模态受损或移除后模型置信度可能反而升高，说明分类精度不能代替概率可靠性\cite{calibratingmultimodallearning2023}。工业异常检测已证明颜色、形态和多视图信息具有互补价值，并形成 RGB/三维基准及混合融合、跨模态映射方法\cite{themvtec3d2022,multimodalindustrialanomaly2023,multimodalindustrialanomaly2024}。然而制造环境相对稳定，难以直接覆盖运输振动、光照变化、封志事件和弱网缺失并存的物流条件；将多模态数据直接拼接或统一填补，仍不能解释各模态在何种扰动下值得信任。

\textbf{（2）多变量时序异常研究能够描述动态关系，但事件级可靠评价与异步原因区分仍不充分。}时空注意、传感器图分解和结构先验可联合刻画变量关系与长程依赖\cite{anomalydetectionfor2024,graphbasedtime2024,fusiongraphstructure2024}；关联差异与动态分解分别缓解异常重建过泛化和非平稳漂移\cite{anomalytransformertime2022,driftdoesnt2023}。这些方法多输出点级异常分数，较少同时处理“物理状态变化、传感质量变化和网络可用性变化”。可靠异常检测基准还指出，随机切窗、点调整和事后阈值可能显著夸大性能，必须采用完整事件划分、事件级指标和冻结阈值\cite{theelephantin2024,deeplearningfor2024}。这要求跨境物流模型不仅识别异常，还应报告不确定性，并在受控模态缺失、时间偏移和噪声矩阵下检验性能退化。

\textbf{（3）边缘智能从计算卸载走向任务相关的信息选择，但通信--计算--能量仍常被分开处理。}边缘计算与边缘智能通过靠近数据源完成推理降低响应时间和带宽开销，相关研究系统讨论了计算卸载、通信/计算资源分配及端侧轻量模型\cite{edgecomputingvision2016,edgeintelligencepaving2019,mobileedgecomputing2017,asurveyon2017,mcunettinydeep2020}。容迟网络利用存储--携带--转发支持间歇连接下的最终交付\cite{adelaytolerant2003}；信息年龄及有用信息年龄用于刻画接收端信息的新鲜度和任务价值\cite{ageofinformation2021,optimizetheage2024}；任务导向和语义通信则推动传输目标由比特重构转向下游任务效用\cite{taskorientedcommunication2022,semanticcommunicationenabling2023}。然而，现有工作往往分别研究卸载、压缩、上传或能耗，难以直接回答物流端侧在电量、缓存、算力和网络同时变化时，何种证据应本地推理、摘要、缓存、续传或优先上传。

\textbf{（4）仓储物联网已能实现多源监测，但复合风险演化与可校准预警仍是薄弱环节。}工业物联网与边云协同已被用于冷库安全监测和无监督异常识别\cite{industrialinternetof2022}，冷链数据研究揭示了概念漂移、稀有异常与类别不平衡问题\cite{ananomalydetection2025}；蜂窝物联网和视觉监测案例也说明传感器、图像与状态告警具有应用基础\cite{deeplearninganomaly2021,intelligentmonitoringand2024}。但温湿度、烟雾、门禁和视频通常由独立规则处理，难以表示缓慢环境漂移与突发安防事件之间的时序关系。仓储研究因而需要由“多源数据接入”推进到“风险状态如何演化、何时需要分级响应以及置信度是否可信”。

\subsubsection{现有研究不足}

综合上述进展，现有研究存在三个相互关联的断点。第一，多模态方法通常假定融合权重或缺失机制相对稳定，尚不能充分解释异步、缺失与噪声共同作用时，振动、倾斜、封志、光感和图像的互补信息何时仍可支持货物异常识别。第二，弱网可靠上传、边缘推理、数据压缩、离线缓存、断点续传和低功耗管理多由不同目标驱动，缺少风险信息价值对通信、计算与能量动作的统一约束；仅减少字节或延长续航，可能以延迟高风险证据为代价。第三，仓储多源监测仍以单传感器阈值或独立分类为主，对不同时标证据的持续性、共现关系和可靠性变化刻画不足，难以同时保证预警提前量、误报控制和概率校准。

相应地，项目需要建立三层可证伪链条：货物多模态可靠表征是否在扰动下保持异常检出与校准；风险价值驱动的边缘策略是否在相同安全约束下改善送达、时延、字节和能耗；仓储多源关系是否真正带来复合异常的提前且可信预警。三项结论均须在完整运输任务或完整仓储事件的外层留出上检验，不能以随机窗口或测试后调阈替代。

\subsubsection{本项目的研究切入点}

针对第一项不足，本项目以振动、倾斜、电子封志、光感和包装图像为对象，联合建模模态内时序、时间偏移、观测质量和跨模态一致性，形成随观测条件变化的可靠性表征，同时输出冲击/倾斜、非法开启和包装破损/变形等事件概率与不确定性。科学假设 H1 是：相对于单模态、静态融合和仅缺失填补模型，该表征应在完整任务外层测试中同时改善事件识别、损伤定位和概率校准，并在预设扰动下表现出更缓退化；否则拒绝该假设。

针对第二项不足，本项目把网络状态、缓存队列、端侧算力、电池和事件风险纳入统一决策状态，依据证据对远端告警风险的边际贡献及按时送达概率，联合选择本地推理深度、压缩级别、缓存优先级、上传/续传和工作周期。科学假设 H2 是：该协同机制能够在不恶化高风险漏报和 P95 告警时延的条件下减少传输字节和单位任务能耗，或在相同资源预算下提高高风险证据按时送达率；若资源收益来自漏报增加或同轨迹比较无稳定优势，则拒绝该假设。

针对第三项不足，本项目融合温湿度、安防视频、烟雾和门禁，建立不同时标、可靠性加权的仓储风险演化模型，并将校准风险映射为分级预警和响应建议。科学假设 H3 是：相对于独立阈值、单源分类和无时序关系模型，跨源关系应提高事件检出率和预警提前量、降低误报并改善校准；若时间或来源置乱后仍有同等收益，或外层仓储事件无稳定增益，则拒绝该假设。运输阶段状态可在货物级配对记录具备时作为可选先验，但仓储主线不以此为前提。

通过上述切入，三条研究内容保持原始指南的完整边界，并由“可靠感知—资源协同—仓储预警”形成连续科学叙事。项目不研发新的蜂窝物理层协议，不以建设综合物流管理平台、优化运输路径或宏观供应链决策为主要目标；现场条件尚未落实时，以受控实验、公开数据和可复现弱网轨迹验证方法边界，不虚构合作、设备和样本规模。

% \subsubsection{科学意义或应用前景}
% % 在此处撰写科学意义或应用前景。
